We determine exactly where the outer sign of a common clustered multiplier can
survive in a polarized two-lane covariance.  In an orthogonal block direct sum,
simultaneous multiplication of both lanes by $C_h$ produces the local factor
$|C_h|^2$; hence every common blockwise unit phase is invisible to covariance
and both lane norms.  For nonorthogonal reassembly and real sign flips, we
derive an exact graph-cut polynomial and prove that invariance under all sign
patterns is equivalent to the vanishing of every symmetrized cross-block edge.
Combining the direct-sum identity with the TPC-243 hard-window bilinear estimate
bounds the difference between any two common sign patterns by
$2\eps\norm W\norm B$.  The literal V59 interpretation remains conditional on
the missing phasewise primitive two-lane coefficient attachment.  Thus the
result is a structural localization theorem and obstruction: the aggregate
outer sign of $C_h$ cannot control the same-block main covariance, while
internal Möbius cancellation in $|C_h|$ remains untouched.
